%=======================================================================
% Chapter 5 — skeleton.
%
% Scope decision (ROADMAP_2026.md §1): this chapter is the ANSYS
% material-law contribution. The JAXFEM material is held back for Keio and
% does NOT appear here.
%
% HOW TO USE THIS FILE
%   Every \TierA block is writable today -- the claim is already true and the
%   evidence is named in the comment above it. Every \TierB block waits on
%   runs (ROADMAP_2026.md §2-3). Prose is deliberately terse: this is a frame
%   with the facts pinned in place, not finished text.
%
%   Copy into the thesis tree and delete the two \newcommand definitions
%   below once the markers have served their purpose.
%
% PORTING (see PORTING.md). A chapter 5 already exists in the thesis
%   repository, so this is material to MERGE INTO it, not a file to drop in
%   whole. It is written to survive that:
%     - packages: amsmath, amssymb, bm, xcolor, graphicx. Nothing else.
%       Numbers are written out rather than via siunitx, and the figure comes
%       in as a PNG, so TikZ is NOT required in the host tree.
%     - every label is namespaced ch5-* so it cannot collide.
%     - one path to set: \ansysfigdir below. Nothing else refers to this repo.
%
% CITATION WARNING, from biofilm_3tooth_refs.bib:
%   growth kinematics F = Fe.Fg, Fg = (1+alpha)I  -> Klempt2024DiffusionDrivenGrowth
%   the 5-species phi/psi interaction dynamics    -> Klempt2026ContinuumBacterialGrowth
%   These are different papers. Citing the 2026 one for the kinematics is the
%   easy mistake here.
%=======================================================================
% ---- the one thing to set when porting -------------------------------
% Directory holding flow_oliver_solution_loop.png (with trailing slash).
\providecommand{\ansysfigdir}{../assets/}
% ----------------------------------------------------------------------

% Drafting markers. Delete both once the chapter is written.
\providecommand{\TierA}[1]{\marginpar{\footnotesize A}#1}   % writable now
\providecommand{\TierB}[1]{\marginpar{\footnotesize B}\textcolor{gray}{#1}} % needs runs

\chapter{A growth constitutive law for an existing FEM framework}
\label{ch:ansys}

%-----------------------------------------------------------------------
\section{Scope and contribution}
\label{sec:ch5-scope}

% TierA. The assignment is THESIS_ASSIGNMENT.md §1; open the chapter from it.
\TierA{%
The task set for this thesis is that the pointwise TMCMC model captures the
temporal evolution at a material point but not its spatial variation, and that
it should be integrated into an existing, working FEM framework so that spatial
effects are represented as well.

This chapter contributes the constitutive half of that integration: a growth
and viscoelastic material law, implemented against the ANSYS user-material
interface, verified against an independent implementation, and packaged so that
it can be called from the partner framework at the point where that framework
already calls its own elastic law. The field solution itself remains with the
partner framework; \S\ref{sec:ch5-integration} makes that division explicit.}

%-----------------------------------------------------------------------
\section{The constitutive model}
\label{sec:ch5-model}

% TierA. Kinematics: cite Klempt2024DiffusionDrivenGrowth, NOT the 2026 paper.
\TierA{%
Growth enters multiplicatively, $\mathbf F = \mathbf F_e \mathbf F_v \mathbf F_g$
with $\mathbf F_g = (1+\alpha)\mathbf I$ \cite{Klempt2024DiffusionDrivenGrowth},
where $\alpha$ is supplied as a field rather than evolved locally. The elastic
response is compressible Mooney--Rivlin, reducing to neo-Hookean at $C_{01}=0$,
and the viscous branch is a deviatoric Newtonian dashpot carrying
$\mathbf F_v$.}

\TierA{%
Written out in the order the implementation evaluates it, a step takes the
total deformation gradient $\mathbf F$ and the viscous state
$\mathbf F_v^{\,n}$ and returns the Cauchy stress and $\mathbf F_v^{\,n+1}$:

\begin{enumerate}
\item the mechanical part removes the growth,
      $\mathbf F_{\mathrm{mech}} = \mathbf F\,\mathbf F_g^{-1}$, and a trial
      elastic gradient follows from the viscous state at the start of the step,
      $\mathbf F_e^{\mathrm{tr}} = \mathbf F_{\mathrm{mech}}(\mathbf F_v^{\,n})^{-1}$;
\item with $\mathbf b_e = \mathbf F_e \mathbf F_e^{\mathsf T}$ and
      $J_e = \det \mathbf F_e$, the deviatoric Kirchhoff stress follows from
      $W = C_{10}(\bar I_1 - 3) + C_{01}(\bar I_2 - 3) + \tfrac{1}{D_1}(J_e-1)^2$,
      which reduces to neo-Hookean at $C_{01}=0$;
\item the viscous state advances by one step of the deviatoric Newtonian flow
      rule,
      \begin{equation}
        \mathbf F_v^{\,n+1} =
        \Bigl( \mathbf I + \frac{\Delta t}{2\,\eta\,J_e}\,
        \bm\tau_{\mathrm{dev}} \Bigr)\mathbf F_v^{\,n},
        \label{eq:ch5-flow}
      \end{equation}
      with $\bm\tau_{\mathrm{dev}}$ taken at the trial state;
\item the elastic gradient and the stress are recomputed at
      $\mathbf F_v^{\,n+1}$, the volumetric term $\tfrac{2}{D_1}(J_e-1)$ is
      added, and the result is returned in the solver's Voigt ordering.
\end{enumerate}

\begin{figure}[htbp]
  \centering
  \includegraphics[width=0.9\linewidth]{\ansysfigdir flow_growth_kinematics.png}
  \caption[Multiplicative growth kinematics]{%
    The three configurations the decomposition introduces: growth carries the
    reference state to an intermediate one that is stress-free but generally
    incompatible, and the elastic part restores compatibility, which is where
    the residual stress comes from. Drawn for $\mathbf F = \mathbf F_e\mathbf F_g$;
    the viscous factor of \S\ref{sec:ch5-model} sits between the two and adds a
    further intermediate configuration without changing the argument.}
  \label{fig:ch5-kinematics}
\end{figure}

Two properties of \eqref{eq:ch5-flow} matter later. It is a single step with no
iteration and no residual: the increment is formed at $\mathbf F_v^{\,n}$ and
applied, so the update is explicit in the flow direction even though the stress
is reported at the end of the step. And $\eta/(2C_{10})$ is a relaxation time
the step has to resolve --- \S\ref{sec:ch5-integration} gives the consequence
and the bound.

The growth driver $\alpha$ is \emph{prescribed}, not evolved: it is supplied
per material point by the ecology model. This chapter is therefore about the
constitutive response to a given growth field, and the verification below
should be read in that light.}

%-----------------------------------------------------------------------
\section{Implementation}
\label{sec:ch5-impl}

% TierA. ansys_usermat/usermat_biofilm.f + README.md.
\TierA{%
The law is implemented as an ANSYS \texttt{USERMAT}.%
\footnote{\texttt{ansys\_usermat/usermat\_biofilm.f} in the accompanying
repository.} Two implementation points are
worth stating because they are where ports of this kind usually fail. First,
the Voigt convention differs between the two solvers used here --- ANSYS orders
the shear components $12,23,13$ and Abaqus $12,13,23$ --- so the two
implementations cannot share an index map. Second, the consistent tangent is
obtained by perturbation of $\mathbf F$ rather than in closed form.}

% TierA. tests/test_tangent_quality.py -- 648 combinations vs exact AD.
\TierA{%
The finite-difference choice invites an obvious objection, so it is worth
answering with a measurement rather than an assurance. A tangent that is
subtly wrong does not produce a wrong answer; it produces a Newton iteration
that fails to converge, and in a handover that failure is very hard for the
recipient to attribute. The tangent was therefore compared against forward-mode
automatic differentiation of the same perturbation --- an independent reference
rather than a second difference quotient --- across the range the routine is
expected to see: stiffness over three decades, Poisson ratios to $0.49$, growth
to $\alpha=0.35$, both material paths, elastic and viscous branches, and time
steps up to the stability limit of \S\ref{sec:ch5-integration}. Over 648
combinations the relative discrepancy has median $1.2\times10^{-7}$ and
worst case $1.1\times10^{-6}$, with no region of the range degrading.
That is a statement about the tangent in isolation; what it does inside a
solver is shown in Figure~\ref{fig:ch5-convergence}.}

%-----------------------------------------------------------------------
\section{Verification}
\label{sec:ch5-verification}

\begin{figure}[htbp]
  \centering
  \includegraphics[width=0.85\linewidth]{\ansysfigdir flow_vv_hierarchy.png}
  \caption[Where this chapter's checks sit]{%
    The verification and validation hierarchy, with this chapter's evidence
    placed on it. Everything in \S\ref{sec:ch5-verification} is \emph{code}
    verification --- does the implementation solve the equations it claims to
    --- and none of it is validation against measurement, which belongs to a
    different chapter and a different kind of evidence.}
  \label{fig:ch5-vv}
\end{figure}

\subsection{Equivalence of two independent implementations}
% TierA. ansys_usermat/crosscheck/README.md -- 0 ULP over 8017 states.
\TierA{%
The same law was implemented twice, independently, against the two solvers'
user-material interfaces. Agreement was then hunted rather than sampled: an
adversarial search over deformation states looked for any state on which the
two disagree. Over 8017 states --- identity and growth, uniaxial, shear,
elastic ($\eta=0$), frozen ($\eta\to\infty$), large growth, both
neo-Hookean and Mooney--Rivlin, and random finite-strain states --- the two
agree to $0$ ULP in stress, viscous state and elastic Jacobian.}

\subsection{Growth kinematics against a closed form}
% TierA. ansys_usermat/apdl/ -- free and constrained growth decks.
\TierA{%
Traction-free isotropic growth admits a closed form: the body expands to
$\mathbf F = (1+\alpha)\mathbf I$ and the stress vanishes identically. The
deck reproduces this, at a largest component of $1.9\times10^{-11}$ against a
constrained answer of $1.0\times10^{-4}$. The case is not redundant beside the
constrained one: full kinematic constraint can mask a sign error in
$\mathbf F_e = \mathbf F\mathbf F_g^{-1}$, because the wrong sign cancels
when the element cannot move at all, and only the free case exposes it.}

% TierA. closed_form_reference.py + tests/test_closed_form_reference.py.
\TierA{%
Fully constrained isotropic growth also admits a closed form, and a sharper one
than it first appears: the elastic deformation is isotropic, so the deviatoric
stress of any potential with an isochoric/volumetric split vanishes and only
$\sigma = (2/D_1)(J_e-1)\mathbf I$ with $J_e=(1+\alpha)^{-3}$ survives ---
independently of the shear constants and of the viscosity, since a deviatoric
flow driver is likewise zero.

This is worth stating separately from the reference values shipped with the
verification decks, which are generated by calling the implementation and
therefore record its behaviour rather than test it. Against the closed form the
implementation differs by exactly the spherical term of
\S\ref{sec:ch5-verification-limits}, at a relative agreement of $10^{-12}$
between the measured discrepancy and its predicted value.}

\TierA{%
That agreement of $10^{-12}$ is between the closed form and the Fortran routine
called directly. Figure~\ref{fig:ch5-closedform} shows the same comparison one
level up, where ANSYS is the one calling: the same element and material, solved
under the two boundary conditions, against the same closed form. Two things are
worth reading off it. The free and constrained answers are separated by seven
decades, which is what the growth kinematics being wired correctly looks like.
And the constrained number the solver returns is reproduced by the ideal
volumetric term \emph{plus} the predicted spurious term of
\S\ref{sec:ch5-verification-limits} --- to every digit MAPDL prints, with that
term accounting for $47\,\%$ of the reported stress. Predicting the
discrepancy rather than absorbing it into a tolerance is what makes this a test
of the analysis and not only of the arithmetic.}

\begin{figure}[htbp]
  \centering
  \includegraphics[width=\linewidth]{\ansysfigdir v222_closed_form.png}
  \caption[The two growth cases, solved in ANSYS]{%
    The growth verification pair as ANSYS 2022~R2 solved it: one SOLID185,
    $\alpha=0.05$, identical material, differing only in the boundary
    conditions. Left: every stress component of both solves. Under free growth
    the closed form is exactly zero and what remains is solver noise, seven
    decades below the constrained answer; the constrained solve puts its three
    normal components on the closed-form line and its shears at machine noise.
    Right: the constrained value decomposed. The ideal continuum term does
    \emph{not} on its own reproduce the solver --- the sum of it and the
    spherical artefact predicted in closed form does. Because the artefact is
    purely spherical it does not enter the von Mises stress this work reports,
    but it is close to half of the hydrostatic number and the figure shows it
    rather than hiding it in a residual.}
  \label{fig:ch5-closedform}
\end{figure}

\TierA{%
Both checks above are made outside any solver: they compare implementations
against each other and against an analytic answer, built with an ordinary
Fortran compiler. That leaves the gap a reviewer should ask about --- whether
the routine behaves the same way once a commercial solver is calling it, with
its own argument conventions, state storage and step control in between.

That gap is closed. The routine was linked into a custom ANSYS 2022~R2
executable and driven through a harness entry point whose only task is to
unpack the state and property arrays and call it, which is the call shape the
partner framework will use. On the fully constrained single element the solver
returns $\sigma_{xx}=\sigma_{yy}=\sigma_{zz} = -1.0193\times10^{-4}$ against a
closed-form $-1.019276\times10^{-4}$, shear components at machine noise, and
the growth variable read back unchanged.

Beyond one element, an already characterised two-layer curved shell of
$12\,240$ elements was re-run unchanged except for the material block. Against
the same geometry solved with the verified core, the delivered routine
reproduces the von Mises field to every digit reported --- minimum
$4.3523\times10^{-9}$, maximum $1.1862\times10^{-5}$, mean
$9.14434842156836\times10^{-6}$ --- on both runs, with no solver errors.

\begin{figure}[htbp]
  \centering
  \includegraphics[width=\linewidth]{\ansysfigdir cylinder_wrapper_3d.png}
  \caption[The curved-shell case, as solved]{%
    The two-layer curved patch the delivered routine was run on, drawn from the
    solver's own listing. Left: the geometry, with the non-growing substrate in
    grey and the growth layer in red --- the layer is on the outside, so it
    hides the substrate from this direction. Centre: a section at mid-length,
    where the interface at $r = 4.3$ and the two node rows through each layer
    are unmistakable; two rows means two elements through the thickness, which
    is why \S\ref{sec:ch5-limitations} treats this mesh as a floor rather than a
    baseline. Right: the distribution of element von Mises stress, bimodal
    because the substrate does not grow. The stress is shown as a distribution
    rather than mapped onto the geometry because the listing carries no element
    connectivity, so there is no honest way to place it in space.}
  \label{fig:ch5-cylinder}
\end{figure}

Two qualifications belong with that second result. It is a consistency check
between two builds rather than a comparison with a closed form, since two
bonded curved layers have no simple analytic answer; and no mesh-convergence
study has been performed on it. It establishes that the repackaging survives
contact with a real solver on a non-trivial mesh. It does not validate the
mesh.}

\TierA{%
That run also settles the question left open in
\S\ref{sec:ch5-impl}. An
inexact tangent does not corrupt the answer --- the residual it is driven
against is the true one --- so it shows up as iteration count, not as error:
the solve takes more equilibrium iterations per substep, and under automatic
stepping the solver bisects to recover. Both are visible in the log, and
Figure~\ref{fig:ch5-convergence} shows what actually happened. Every substep
reaches the convergence criterion in between one and four equilibrium
iterations, and \texttt{AUTOTS}, free to bisect throughout, instead enlarged
the increment at every opportunity, from $0.5$ to $1.6875$; the load step
completed in six substeps and thirteen cumulative iterations with no cut-back.
The finite-difference tangent is therefore not merely accurate in isolation but
accurate enough that the solver never has to work around it.}

\begin{figure}[htbp]
  \centering
  \includegraphics[width=\linewidth]{\ansysfigdir v222_convergence.png}
  \caption[Newton convergence with the finite-difference tangent]{%
    Convergence of the $12\,240$-element solve with the delivered routine, read
    from the solver's own listing. Left: force residual against equilibrium
    iteration, one curve per substep, with the ANSYS criterion dashed and
    converged states ringed. The fall is not strictly monotone --- substep~1
    rises once between its second and third iterations --- which is ordinary
    under an inexact tangent and is shown rather than smoothed. Right: the time
    increment the solver chose for each substep, annotated with the iterations
    that substep needed. The increment grows monotonically; the short final
    substep is the remainder to the end of the load step, not a cut-back. The
    single-element case is deliberately absent: it is fully constrained, so it
    has no free degrees of freedom and its residual is identically zero in both
    builds, which would look like agreement while carrying no information.}
  \label{fig:ch5-convergence}
\end{figure}

\subsection{What these verifications establish, and what they do not}
\label{sec:ch5-verification-limits}
% TierA. DEVIATOR_SCALING_FINDING.md. This subsection is the most valuable
% one in the chapter and the one most likely to be asked about in the viva.
\TierA{%
Cross-implementation agreement establishes that two ports of the same
expression compute that expression identically. It does not establish that the
expression is right, and the distinction is not hypothetical here.

The isochoric split in both implementations applies $J^{-2/3}$ to the
subtracted trace but not to the tensor beside it. The coded and correct forms
coincide exactly at $J=1$ and nowhere else --- and growth, by construction,
moves $J_e$ away from $1$. The resulting discrepancy is
\emph{purely spherical} for any deformation, so it is a pressure error: the
deviatoric stress, and hence the von Mises stress reported throughout this
work, is unaffected to machine precision. Its indirect route --- a spherical
term in the viscous flow driver, which makes $\det\mathbf F_v$ drift from
unity --- was measured at $\le 0.01\,\%$ in von Mises over 40 steps.

\begin{figure}[htbp]
  \centering
  \includegraphics[width=\linewidth]{\ansysfigdir flow_vv_blindspot.png}
  \caption[What the checks cover, and what they miss]{%
    Four checks that look independent, and the one error none of them can see.
    In each case the reason is structural rather than a matter of tolerance:
    two ports of the same expression, a case that lands exactly where both
    forms agree, reference values generated by the implementation itself, and a
    metric that cannot represent a spherical term.}
  \label{fig:ch5-blindspot}
\end{figure}

The point worth drawing out is why none of the three checks above sees it. The
cross-implementation hunt compares two ports carrying the same expression. The
free-growth deck lands on $J_e=1$ exactly, the single case where both forms
agree. And the constrained-growth reference values are generated by calling the
implementation, so they record its behaviour rather than test it. A fourth
check, a growth cross-comparison against a second material model, compares von
Mises and is therefore blind to a spherical term by construction.}

\TierA{%
One route is left unmeasured, and is stated as such. The discrepancy is a
pressure error, and in an equilibrium solve a wrong pressure changes the
displacement field, which reaches the von Mises stress indirectly even though
it cannot touch it at a fixed deformation gradient. Bounding that would mean
solving the same problem both ways. The one place it has been looked at ---
two routines with different elastic potentials agreeing to $0.1\,\%$ on
constrained growth --- was at $\nu = 0.49$, the most forgiving case for a
pressure error, since a nearly incompressible material resists the volumetric
strain such an error would produce. Lower Poisson ratios are untested.}

%-----------------------------------------------------------------------
\section{Integration into the partner framework}
\label{sec:ch5-integration}

\subsection{The framework}
% TierA. Figure already built: assets/flow_oliver_solution_loop.png
\TierA{%
The partner framework solves the mechanics in ANSYS and the transport fields in
its own user-programmable routines, alternating substep by substep
(Figure~\ref{fig:partner-loop}). The transport fields are not ANSYS degrees of
freedom; they live in the framework's own data pool, assembled on a meshless
neighbour operator and solved with a sparse direct solver.}

\begin{figure}[htbp]
  \centering
  \includegraphics[width=\linewidth]{\ansysfigdir flow_oliver_solution_loop.png}
  % Vector alternative, if the host tree already loads TikZ with
  % shapes.geometric, arrows.meta, positioning, fit, backgrounds, calc:
  %   \resizebox{\linewidth}{!}{% flow_oliver_solution_loop.tex — the solution loop of the partner ANSYS/UPF
% framework, and where our constitutive law attaches to it.
%
% Traced from the source pool (USolBeg / Usermat / USSFin), as recorded in
% ansys_usermat/OLIVER_MODEL_NOTES.md §"How the whole thing runs". Everything
% in solid boxes is read directly from the code; the one interpretation is
% marked as such. Usage:
%   \resizebox{\linewidth}{!}{% flow_oliver_solution_loop.tex — the solution loop of the partner ANSYS/UPF
% framework, and where our constitutive law attaches to it.
%
% Traced from the source pool (USolBeg / Usermat / USSFin), as recorded in
% ansys_usermat/OLIVER_MODEL_NOTES.md §"How the whole thing runs". Everything
% in solid boxes is read directly from the code; the one interpretation is
% marked as such. Usage:
%   \resizebox{\linewidth}{!}{% flow_oliver_solution_loop.tex — the solution loop of the partner ANSYS/UPF
% framework, and where our constitutive law attaches to it.
%
% Traced from the source pool (USolBeg / Usermat / USSFin), as recorded in
% ansys_usermat/OLIVER_MODEL_NOTES.md §"How the whole thing runs". Everything
% in solid boxes is read directly from the code; the one interpretation is
% marked as such. Usage:
%   \resizebox{\linewidth}{!}{\input{flow_oliver_solution_loop.tex}}
\begin{tikzpicture}[
  every node/.style = {font=\small},
  N/.style    = {draw, rounded corners=3pt, align=left, thick,
                 text width=9.6cm, inner sep=6pt},
  Nonce/.style= {N, fill=gray!10},
  Ngp/.style  = {N, fill=orange!12},
  Nfld/.style = {N, fill=teal!12},
  Nours/.style= {draw, rounded corners=3pt, align=left, very thick,
                 draw=red!65!black, fill=red!8, text width=6.4cm, inner sep=6pt},
  ar/.style   = {-Stealth, thick},
  arb/.style  = {-Stealth, thick, dashed, gray!70},
  lbl/.style  = {font=\scriptsize\itshape, fill=white, inner sep=1.5pt, align=center},
  note/.style = {font=\scriptsize\itshape, align=left, text width=6.2cm, gray!45!black},
]

% ── entry ───────────────────────────────────────────────────────────
\node[Nonce] (top) {%
  \textbf{ANSYS solve} — \texttt{SOLID185}, \texttt{NLGEOM,ON}\\
  custom \texttt{ANSYS.exe}, UPF routines linked in};

% ── USolBeg ─────────────────────────────────────────────────────────
\node[Nonce, below=0.55cm of top] (beg) {%
  \textbf{\texttt{USolBeg}} \hfill \textit{once, at solution start}\\[4pt]
  $\sim$150 \texttt{parevl} calls: every APDL parameter $\rightarrow$
  \texttt{/usercm/} common block\\[2pt]
  \texttt{InitVals} — allocate the shared (MPI) data arena\\[2pt]
  \texttt{NEM\_CreateData\_Init} — build the meshless neighbour operator\\[2pt]
  \texttt{Mapping\_Node\_ID}, \texttt{Initial\_Values}};
\draw[ar] (top) -- (beg);

% ── the staggered loop ──────────────────────────────────────────────
\node[Ngp, below=0.95cm of beg] (gp) {%
  \textbf{\texttt{usermat}} \hfill \textit{per Gauss point, per equilibrium iteration}\\[4pt]
  \texttt{GetVals} / \texttt{GetTMP} — read this point's state from the pool\\[2pt]
  \textbf{\texttt{CALL AceGenNeoHookV04}} $\rightarrow$ $\bm\sigma$, \texttt{dsdePl}\\[2pt]
  \texttt{SetVals} — write state back};
\draw[ar] (beg) -- node[lbl,right,xshift=2pt] {mechanics} (gp);

\node[Nfld, below=0.95cm of gp] (fld) {%
  \textbf{\texttt{USSFin}} \hfill \textit{after each substep}\\[4pt]
  assemble $+$ PARDISO solve $\rightarrow$ temperature field\\[2pt]
  assemble $+$ PARDISO solve $\rightarrow$ \texttt{Nut1} field\\[2pt]
  assemble $+$ PARDISO solve $\rightarrow$ \texttt{Nut2} field\\[2pt]
  \textcolor{gray!55!black}{\texttt{AGPhaseViskoP21V07} — commented out}\\[2pt]
  \textcolor{gray!55!black}{\texttt{CalcLaserIntegralOMP} / \texttt{CALCPYRO}
    — laser \& pyrometer (inherited glass process)}};
\draw[ar] (gp) -- node[lbl,right,xshift=2pt] {transport} (fld);

% staggered feedback
\draw[arb] (fld.west) -- ++(-0.85,0) |- node[lbl,pos=0.25,left,xshift=-2pt]
  {next substep} (gp.west);

% ── our attachment point ────────────────────────────────────────────
\node[Nours, right=1.15cm of gp] (ours) {%
  \textbf{our contribution attaches here}\\[4pt]
  \texttt{BIOFILM\_GROWTH\_VISCO\_V01}\\
  \scriptsize(\texttt{ansys\_usermat/biofilm\_material\_v01.f})\\[4pt]
  \scriptsize same signature shape as \texttt{AceGenNeoHookV04}, plus the
  growth variable $\alpha$ and the viscous state $\mathbf F_v$\\[2pt]
  \scriptsize adds $\mathbf F=\mathbf F_e\mathbf F_v\mathbf F_g$,
  $\mathbf F_g=(1{+}\alpha)\mathbf I$ over their purely elastic law};
\draw[ar, draw=red!65!black, very thick] (ours) -- (gp);

% ── annotations ─────────────────────────────────────────────────────
\node[note, right=1.15cm of fld, yshift=0.1cm] (n1) {%
  Staggered / operator-split: ANSYS solves the mechanics, \texttt{USSFin}
  solves the transport fields on the NEM operator with MKL PARDISO, and the
  two alternate substep by substep.\\[3pt]
  The transport fields are \textbf{not} ANSYS degrees of freedom --- they live
  entirely in the UPF's own data pool.\\[3pt]
  \textcolor{red!55!black}{Interpretation, not stated by the code:} the two
  \texttt{Nut} systems are most likely the biofilm and nutrient fields rather
  than two nutrients.};

\node[note, right=1.15cm of beg, yshift=0.05cm] {%
  Because every parameter arrives through \texttt{parevl} into
  \texttt{/usercm/}, the ANSYS \texttt{prop()} array is unused --- material
  constants do \emph{not} come in through \texttt{TB,USER} here, unlike our
  own deck.};

\end{tikzpicture}
}
\begin{tikzpicture}[
  every node/.style = {font=\small},
  N/.style    = {draw, rounded corners=3pt, align=left, thick,
                 text width=9.6cm, inner sep=6pt},
  Nonce/.style= {N, fill=gray!10},
  Ngp/.style  = {N, fill=orange!12},
  Nfld/.style = {N, fill=teal!12},
  Nours/.style= {draw, rounded corners=3pt, align=left, very thick,
                 draw=red!65!black, fill=red!8, text width=6.4cm, inner sep=6pt},
  ar/.style   = {-Stealth, thick},
  arb/.style  = {-Stealth, thick, dashed, gray!70},
  lbl/.style  = {font=\scriptsize\itshape, fill=white, inner sep=1.5pt, align=center},
  note/.style = {font=\scriptsize\itshape, align=left, text width=6.2cm, gray!45!black},
]

% ── entry ───────────────────────────────────────────────────────────
\node[Nonce] (top) {%
  \textbf{ANSYS solve} — \texttt{SOLID185}, \texttt{NLGEOM,ON}\\
  custom \texttt{ANSYS.exe}, UPF routines linked in};

% ── USolBeg ─────────────────────────────────────────────────────────
\node[Nonce, below=0.55cm of top] (beg) {%
  \textbf{\texttt{USolBeg}} \hfill \textit{once, at solution start}\\[4pt]
  $\sim$150 \texttt{parevl} calls: every APDL parameter $\rightarrow$
  \texttt{/usercm/} common block\\[2pt]
  \texttt{InitVals} — allocate the shared (MPI) data arena\\[2pt]
  \texttt{NEM\_CreateData\_Init} — build the meshless neighbour operator\\[2pt]
  \texttt{Mapping\_Node\_ID}, \texttt{Initial\_Values}};
\draw[ar] (top) -- (beg);

% ── the staggered loop ──────────────────────────────────────────────
\node[Ngp, below=0.95cm of beg] (gp) {%
  \textbf{\texttt{usermat}} \hfill \textit{per Gauss point, per equilibrium iteration}\\[4pt]
  \texttt{GetVals} / \texttt{GetTMP} — read this point's state from the pool\\[2pt]
  \textbf{\texttt{CALL AceGenNeoHookV04}} $\rightarrow$ $\bm\sigma$, \texttt{dsdePl}\\[2pt]
  \texttt{SetVals} — write state back};
\draw[ar] (beg) -- node[lbl,right,xshift=2pt] {mechanics} (gp);

\node[Nfld, below=0.95cm of gp] (fld) {%
  \textbf{\texttt{USSFin}} \hfill \textit{after each substep}\\[4pt]
  assemble $+$ PARDISO solve $\rightarrow$ temperature field\\[2pt]
  assemble $+$ PARDISO solve $\rightarrow$ \texttt{Nut1} field\\[2pt]
  assemble $+$ PARDISO solve $\rightarrow$ \texttt{Nut2} field\\[2pt]
  \textcolor{gray!55!black}{\texttt{AGPhaseViskoP21V07} — commented out}\\[2pt]
  \textcolor{gray!55!black}{\texttt{CalcLaserIntegralOMP} / \texttt{CALCPYRO}
    — laser \& pyrometer (inherited glass process)}};
\draw[ar] (gp) -- node[lbl,right,xshift=2pt] {transport} (fld);

% staggered feedback
\draw[arb] (fld.west) -- ++(-0.85,0) |- node[lbl,pos=0.25,left,xshift=-2pt]
  {next substep} (gp.west);

% ── our attachment point ────────────────────────────────────────────
\node[Nours, right=1.15cm of gp] (ours) {%
  \textbf{our contribution attaches here}\\[4pt]
  \texttt{BIOFILM\_GROWTH\_VISCO\_V01}\\
  \scriptsize(\texttt{ansys\_usermat/biofilm\_material\_v01.f})\\[4pt]
  \scriptsize same signature shape as \texttt{AceGenNeoHookV04}, plus the
  growth variable $\alpha$ and the viscous state $\mathbf F_v$\\[2pt]
  \scriptsize adds $\mathbf F=\mathbf F_e\mathbf F_v\mathbf F_g$,
  $\mathbf F_g=(1{+}\alpha)\mathbf I$ over their purely elastic law};
\draw[ar, draw=red!65!black, very thick] (ours) -- (gp);

% ── annotations ─────────────────────────────────────────────────────
\node[note, right=1.15cm of fld, yshift=0.1cm] (n1) {%
  Staggered / operator-split: ANSYS solves the mechanics, \texttt{USSFin}
  solves the transport fields on the NEM operator with MKL PARDISO, and the
  two alternate substep by substep.\\[3pt]
  The transport fields are \textbf{not} ANSYS degrees of freedom --- they live
  entirely in the UPF's own data pool.\\[3pt]
  \textcolor{red!55!black}{Interpretation, not stated by the code:} the two
  \texttt{Nut} systems are most likely the biofilm and nutrient fields rather
  than two nutrients.};

\node[note, right=1.15cm of beg, yshift=0.05cm] {%
  Because every parameter arrives through \texttt{parevl} into
  \texttt{/usercm/}, the ANSYS \texttt{prop()} array is unused --- material
  constants do \emph{not} come in through \texttt{TB,USER} here, unlike our
  own deck.};

\end{tikzpicture}
}
\begin{tikzpicture}[
  every node/.style = {font=\small},
  N/.style    = {draw, rounded corners=3pt, align=left, thick,
                 text width=9.6cm, inner sep=6pt},
  Nonce/.style= {N, fill=gray!10},
  Ngp/.style  = {N, fill=orange!12},
  Nfld/.style = {N, fill=teal!12},
  Nours/.style= {draw, rounded corners=3pt, align=left, very thick,
                 draw=red!65!black, fill=red!8, text width=6.4cm, inner sep=6pt},
  ar/.style   = {-Stealth, thick},
  arb/.style  = {-Stealth, thick, dashed, gray!70},
  lbl/.style  = {font=\scriptsize\itshape, fill=white, inner sep=1.5pt, align=center},
  note/.style = {font=\scriptsize\itshape, align=left, text width=6.2cm, gray!45!black},
]

% ── entry ───────────────────────────────────────────────────────────
\node[Nonce] (top) {%
  \textbf{ANSYS solve} — \texttt{SOLID185}, \texttt{NLGEOM,ON}\\
  custom \texttt{ANSYS.exe}, UPF routines linked in};

% ── USolBeg ─────────────────────────────────────────────────────────
\node[Nonce, below=0.55cm of top] (beg) {%
  \textbf{\texttt{USolBeg}} \hfill \textit{once, at solution start}\\[4pt]
  $\sim$150 \texttt{parevl} calls: every APDL parameter $\rightarrow$
  \texttt{/usercm/} common block\\[2pt]
  \texttt{InitVals} — allocate the shared (MPI) data arena\\[2pt]
  \texttt{NEM\_CreateData\_Init} — build the meshless neighbour operator\\[2pt]
  \texttt{Mapping\_Node\_ID}, \texttt{Initial\_Values}};
\draw[ar] (top) -- (beg);

% ── the staggered loop ──────────────────────────────────────────────
\node[Ngp, below=0.95cm of beg] (gp) {%
  \textbf{\texttt{usermat}} \hfill \textit{per Gauss point, per equilibrium iteration}\\[4pt]
  \texttt{GetVals} / \texttt{GetTMP} — read this point's state from the pool\\[2pt]
  \textbf{\texttt{CALL AceGenNeoHookV04}} $\rightarrow$ $\bm\sigma$, \texttt{dsdePl}\\[2pt]
  \texttt{SetVals} — write state back};
\draw[ar] (beg) -- node[lbl,right,xshift=2pt] {mechanics} (gp);

\node[Nfld, below=0.95cm of gp] (fld) {%
  \textbf{\texttt{USSFin}} \hfill \textit{after each substep}\\[4pt]
  assemble $+$ PARDISO solve $\rightarrow$ temperature field\\[2pt]
  assemble $+$ PARDISO solve $\rightarrow$ \texttt{Nut1} field\\[2pt]
  assemble $+$ PARDISO solve $\rightarrow$ \texttt{Nut2} field\\[2pt]
  \textcolor{gray!55!black}{\texttt{AGPhaseViskoP21V07} — commented out}\\[2pt]
  \textcolor{gray!55!black}{\texttt{CalcLaserIntegralOMP} / \texttt{CALCPYRO}
    — laser \& pyrometer (inherited glass process)}};
\draw[ar] (gp) -- node[lbl,right,xshift=2pt] {transport} (fld);

% staggered feedback
\draw[arb] (fld.west) -- ++(-0.85,0) |- node[lbl,pos=0.25,left,xshift=-2pt]
  {next substep} (gp.west);

% ── our attachment point ────────────────────────────────────────────
\node[Nours, right=1.15cm of gp] (ours) {%
  \textbf{our contribution attaches here}\\[4pt]
  \texttt{BIOFILM\_GROWTH\_VISCO\_V01}\\
  \scriptsize(\texttt{ansys\_usermat/biofilm\_material\_v01.f})\\[4pt]
  \scriptsize same signature shape as \texttt{AceGenNeoHookV04}, plus the
  growth variable $\alpha$ and the viscous state $\mathbf F_v$\\[2pt]
  \scriptsize adds $\mathbf F=\mathbf F_e\mathbf F_v\mathbf F_g$,
  $\mathbf F_g=(1{+}\alpha)\mathbf I$ over their purely elastic law};
\draw[ar, draw=red!65!black, very thick] (ours) -- (gp);

% ── annotations ─────────────────────────────────────────────────────
\node[note, right=1.15cm of fld, yshift=0.1cm] (n1) {%
  Staggered / operator-split: ANSYS solves the mechanics, \texttt{USSFin}
  solves the transport fields on the NEM operator with MKL PARDISO, and the
  two alternate substep by substep.\\[3pt]
  The transport fields are \textbf{not} ANSYS degrees of freedom --- they live
  entirely in the UPF's own data pool.\\[3pt]
  \textcolor{red!55!black}{Interpretation, not stated by the code:} the two
  \texttt{Nut} systems are most likely the biofilm and nutrient fields rather
  than two nutrients.};

\node[note, right=1.15cm of beg, yshift=0.05cm] {%
  Because every parameter arrives through \texttt{parevl} into
  \texttt{/usercm/}, the ANSYS \texttt{prop()} array is unused --- material
  constants do \emph{not} come in through \texttt{TB,USER} here, unlike our
  own deck.};

\end{tikzpicture}
}
  % Japanese version: flow_oliver_solution_loop_ja.png (uplatex only).
  \caption[Partner framework solution loop]{%
    Solution loop of the partner framework, traced from source, and the point
    at which the law of this chapter attaches. Solid boxes are read directly
    from the code; the single inference is marked as such in the figure.}
  \label{fig:partner-loop}
\end{figure}

\subsection{The delivered routine}
% TierA. ansys_usermat/biofilm_material_v01.f + tests/test_material_wrapper.py
\TierA{%
The law is delivered as a routine matching the shape of the framework's own
constitutive calls, taking per-phase $(E,\nu)$ blended by a biofilm fraction
and returning Cauchy stress and consistent tangent. Like the routine it
replaces it uses no solver includes, no common block and no
user-programmable-feature calls, so it is independent of the ANSYS release it
is linked into; only the surrounding glue is release-specific.

It is deliberately an adapter around the verified core rather than a
reimplementation, and this is tested rather than asserted: fed $(E,\nu)$, the
delivered routine must reproduce the core fed the $(C_{10},C_{01},D_1)$ those
map to, at zero tolerance. It does, so the equivalence of
\S\ref{sec:ch5-verification} carries over to the delivered routine unchanged.}

\subsection{Interface constraints}
% TierA. Both measured, not assumed. biofilm_material_v01.f header notes 2-3.
\TierA{%
Two constraints were established by measurement and must travel with the
routine, because the caller controls both.

The viscous state must be stored as a full $\mathbf F_v$, not as a
six-component Cauchy--Green tensor: $\mathbf F_v$ does not remain symmetric
under this update, measured at $\sim 6\times10^{-5}$ relative asymmetry after 20
steps over 200 random states, so nine state slots are required where six might
be expected.

The time step must resolve the viscous relaxation time $\eta/(2C_{10})$. The
flow increment is evaluated at the old state, so accuracy degrades as
$\Delta t$ approaches that time; past $\Delta t/\tau \approx 0.5$ the stress
changes sign and past $\Delta t/\tau \approx 1$ it diverges. Setting $\eta=0$
selects the elastic path and removes the constraint.}

% TierA. ansys_usermat/apdl/plot_state_roundtrip.py -- SVAR read-back.
\TierA{%
Those nine slots, and the tenth carrying $\alpha$, are the whole growth
mechanism in ANSYS: with no \texttt{INISTATE} available for a
\texttt{USERMAT}, $\alpha$ cannot be supplied as a field and is written per
material with \texttt{TBDATA} instead. That makes the state array a single
point of failure, and one that fails silently in both of its modes ---
\texttt{TBDATA} accepts at most six values per call and discards the rest
without an error, and \texttt{TB,STATE} declared too short has the same
effect. Either leaves $\alpha$ at zero, which is a completely successful,
completely elastic run that reports no problem. Reading the slots back is the
only direct evidence the write survived, and
Figure~\ref{fig:ch5-stateroundtrip} does so for the three runs that printed
them.}

\begin{figure}[htbp]
  \centering
  \includegraphics[width=\linewidth]{\ansysfigdir v222_state_roundtrip.png}
  \caption[State slots read back out of the solve]{%
    Every \texttt{TB,STATE} slot the three v222 runs printed, compared with
    what the deck wrote. All $176\,800$ samples agree to the five digits MAPDL
    lists. Two limits are marked on the figure rather than left to be assumed:
    $\eta=0$ in all three decks, so $\mathbf F_v$ has nothing to evolve and
    what is shown is the round trip rather than the viscous update; and the
    curved shell lists state under \texttt{ESEL,S,MAT,,2}, so it evidences the
    growth layer alone and the substrate's $\alpha=0$ is not read back
    anywhere. That selection is itself informative --- an earlier version of
    the deck put both volumes on the default material, where the same
    \texttt{ESEL} selected no elements at all, so the $11\,040$ it selects
    here is the check that the two-material assignment took.}
  \label{fig:ch5-stateroundtrip}
\end{figure}

%-----------------------------------------------------------------------
\section{Results}
\label{sec:ch5-results}

\TierB{%
[Awaiting runs. Comparison of von Mises fields across conditions, produced with
the law of this chapter. Reachable through the local v222 build alone --- it
does not depend on the partner framework wiring the routine in
(ROADMAP\_2026.md §3).]}

% TODO(runs): before reporting anything here, confirm dt is inside the stable
% range of §5.5.3 for the actual job. This is cheap to check and expensive to
% discover afterwards.

%-----------------------------------------------------------------------
\section{Limitations}
\label{sec:ch5-limitations}

\TierA{%
The growth driver $\alpha$ is prescribed rather than solved here; the
verification of \S\ref{sec:ch5-verification} therefore covers the constitutive
response to a given field, not the field itself.}

\TierA{%
The pressure discrepancy of \S\ref{sec:ch5-verification-limits} is documented
rather than corrected. Correcting it would invalidate every reference value in
the verification chain for a change that does not move the reported von Mises
results, and it is not yet established whether the published formulation
defines the split as implemented --- in which case the change would be a
disagreement with that formulation rather than a correction.}

\TierA{%
Two further limitations bound what the comparison in
\S\ref{sec:ch5-results} can be read to mean.

The first is a step restriction rather than a modelling choice. Because the
viscous increment of \eqref{eq:ch5-flow} is formed at the previous state, the
time step must resolve $\eta/(2C_{10})$. The routine refuses a step past half
that time, but accuracy degrades continuously below the refusal threshold, so
an accepted step is not thereby an accurate one. Any reported stress is
conditional on the step used, which is recorded with the results rather than
left implicit.

The second concerns the material constants rather than the mechanics. The
composition-to-stiffness bridge that supplies $(E,\nu)$ clips its output to a
physical range, and at the present calibration that bound is reached over much
of the commensal half of composition space rather than at its edge. Two
compositions the model separates can therefore enter the mechanics with the
same stiffness and leave it with the same stress. Where that happens, a
reported difference between conditions is a lower bound on the difference the
model predicts, not an estimate of it.

Sensitivity to the species-stiffness assignment is reported elsewhere as a band
rather than a point value and is not repeated here.}
